\chapter{The ramified Miura parameter}
\label{rpd:chapter}

The Miura $W_4$ algebra has central charge $3-60a^2$.
The simultaneous substitutions $a\mapsto-a$ and $a_i\mapsto-a_i$
for its four Miura currents preserve the central charge and change
the sign of the state of weight three. At $a=0$ this symmetry fixes
the parameter and still
acts nontrivially on the states. A family indexed only by central
charge must therefore be distinguished from the family with this
specified symmetry.

Put $A=\mathbb C[t]$, $R=\mathbb C[a]$, with $t\mapsto a^2$, and work on
$X_R=\mathbb A^1_R$ with coordinate $z$, translation $T=\partial_z$, and
density $dz$. The parameters have cohomological degree and weight zero,
and $T$ raises weight by one. Tensor products in this chapter are tensor products of
field state modules unless stated otherwise. They do not replace cut
convolution or the diagonal tensor of geometric coefficient systems.
In particular, their vertex vacuum does not identify the scalar cut
unit $K$ with the local geometric vacuum $J$.

\section{The polynomial family and its symmetry}

We use the Miura fields $U_0=\mathbf1$, $U_2,U_3,U_4$ and the curved dual
of Theorem~\ref{w4d:thm:main}. Write
\[
 \Gamma_0=S,\qquad \Gamma_2=\gamma_2,\quad
 \Gamma_3=\gamma_3,\quad\Gamma_4=\gamma_4.
\]
Every $U_i$ is even of cohomological degree zero. Every $\Gamma_i$ is odd
of degree one, with weight $1-i$. The scalar-dual state $S$ is not the
vertex vacuum. Modes are normalized by
$Y(x,u)y=\sum_n x_{(n)}y\,u^{-n-1}$. In particular, the normal product is
the $(-1)$-product. The curvature convention is
\[
 d=S_{(0)}|_B,\qquad w=\tfrac12S_{(0)}S,
 \qquad d^2=w_{(0)},\qquad dw\in TB.
\]
The scalar-extended generator bundle $N_R^\circ$ is free of rank four
on $X_R$, with frame $e_0,e_2,e_3,e_4$. Its duality is contravariant:
\[
 \check N_R^\circ=s^{-1}\mathcal{H}om_{\mathcal O_{X_R}}(N_R^\circ,\mathcal O_{X_R})
       \otimes\omega_{X_R}^{-1}.
\]
The canonical line here is relative to $R$:
$\omega_{X_R}^{-1}$ denotes $\omega_{X_R/R}^{-1}$.
The shift $s^{-1}$ places the dual generators in cohomological degree one.
On $X_R^2=X_R\times_{\operatorname{Spec}R}X_R$, the external
two-generator pairing takes values in $s^{-2}\omega_{X_R^2/R}^{-1}$. The local pole coordinate is
$u=z_1-z_2$; its sign is fixed, and the relations are extended from
$u\ne0$ across that diagonal by the stated lattices. No real boundary
or orientation-reversing operation on the curve is used. This finite
locally free duality is not the continuous dual of the full ordered
coefficient coalgebra.

Give $U_s$ weight and filtration degree $s$.
For the odd generators use the positive filtration
\[
 \delta(S)=9,\qquad\delta(\gamma_2)=5,\qquad
 \delta(\gamma_3)=3,\qquad\delta(\gamma_4)=1,
 \qquad\delta(T)=0.
\]
The full coupling and its induced operator are
\[
 q=\sum_{s=2}^4U_s\otimes\gamma_s,\qquad
 e=\mathbf1\otimes S+q,\qquad
 Q=e_{(0)}|_{V\otimes B}=d+q_{(0)}.
\]
One has $e_{(n)}e=0$ for all $n\geq0$, and consequently $Q^2=0$,
by Proposition~\ref{w4d:prop:coupling}. The exact state $e$, including
its derivatives, is part of the data. No augmentation to $R$ is specified.

\begin{lemma}[Polynomial PBW extension]\label{rpd:lem:polynomial}
The universal Miura datum and its curved field dual admit free
$R$-forms $V_R$, $B_R\subset B_R^\circ$ with the full generic ordered PBW
bases. All modes, translation, the inclusion, $S,d,w$, and $e$ preserve
these forms. For every commutative $R$-algebra $D$, their scalar
extensions are the universal field algebras of the specialized
presentations. This includes nonflat $D$.
\end{lemma}
\begin{proof}
Take the $R$-spans of the ordered PBW words inside the generic algebras
over $\mathbb C(a)$, using Propositions~\ref{w4d:prop:lattice}
and~\ref{w4d:prop:dual-PBW}. The coefficient tables
\ref{w4d:tab:C} and~\ref{w4d:tab:E} have entries in $\mathbb Q[a]$.
The primal relation matrix $M$ has diagonal entries one and satisfies
$(1-M)^4=0$. Its inverse is therefore
$M^{-1}=\sum_{j=0}^3(1-M)^j$. The annihilator matrix
$D_M=M^{-\mathrm t}$ therefore has entries in $R[u^{-1}]$,
and $MD_M^{\mathrm t}=1$ holds over that ring.

Normal ordering of a fixed product uses finitely many reductions.
Every dual correction has $m=k+l-i-j\geq2$ and lowers the
filtration by $2m$. A primal correction lowers the filtration by
$m=i+j-k-l\geq2$. Wrongly ordered
creation modes are exchanged using the commutator, and a repeated odd
creation mode is half its commutator. The only numerical divisions are
by nonzero rational integers. The reductions therefore have
coefficients in $R$ and terminate for each pair of finite states.
All modes preserve the indicated $R$-spans.

The generic PBW words are linearly independent, so these spans are
free. Every vertex identity is an identity between elements of these
free modules and holds after extension to $\mathbb C(a)$; hence it holds
over $R$. The same argument applies to the curvature and coupling
identities. After extension to any $D$, the normal-ordering relations
still span the ordered words. The free base-changed model gives their
linear independence, so the universal presentation maps isomorphically
to that model. No flatness of $D$ is used.
The finite relation lattices and their annihilators also commute with
this base change: the specialized matrices still give dual free frames,
by $MD_M^{\mathrm t}=1$. This argument uses no interchange with an
infinite pro-limit.\qedhere
\end{proof}

This assertion concerns universal PBW algebras. It does not identify a
specialization with a simple quotient or with the image of a
specialized free-field representation. It also does not construct the
completed geometric bar comparison.

\begin{lemma}[Signed symmetry]\label{rpd:lem:symmetry}
There is a semilinear involution $\sigma$ of these three algebras
such that
\[
 \sigma(a)=-a,\qquad \sigma(U_s)=(-1)^sU_s,
 \qquad \sigma(\Gamma_s)=(-1)^s\Gamma_s.
\]
It commutes with translation, fixes the vacuum and $S,d,w,e$, and
preserves every mode and the inclusion $B_R\subset B_R^\circ$.
\end{lemma}
\begin{proof}
On the Miura expression the simultaneous change $a\mapsto-a$ and
$a_i\mapsto-a_i$ multiplies each of the four factors by $-1$.
Comparing the coefficients of $(aT)^{4-s}$ gives the stated primal
sign. For the relation matrices, indexed by ordered pairs from
$I=\{0,2,3,4\}$, put
\[
 P_{ij,ij}=(-1)^{i+j}.
\]
Their coefficient identities are
\[
 M(-a)=PM(a)P,\qquad D_M(-a)=PD_M(a)P.
\]
Indeed each coefficient from row $ij$ to column $kl$ has parameter
parity $i+j+k+l$; this is visible term by term in both polynomial
tables. Thus the universal relations are preserved. The maps extend
to all states, and applying them twice is the identity. Since $S$ is
fixed, its zero mode and half-self-product are fixed. Each summand
$U_s\otimes\gamma_s$ has two equal signs, so $e$ is fixed.\qedhere
\end{proof}

The sign in this lemma is not ghost parity. For example,
$\gamma_2$ and $\gamma_4$ are odd states fixed by $\sigma$.

\section{Descent and the fixed fibre}

\begin{definition}\label{rpd:def:system}
A \emph{marked field system} over a commutative $\mathbb C$-algebra $D$
consists of unital graded vertex superalgebras $V$ and
$B\subset B^\circ$, their translations and weight gradings, positive
generator filtrations, a degree-one state $S\in B^\circ$, and a
degree-one state $e\in V\otimes_D B^\circ$.
Require $S_{(n)}B\subset B$ and $S_{(n)}S\in B$ for $n\geq0$,
the curved identities for $d=S_{(0)}|_B$ and
$w=S_{(0)}S/2$, and $e_{(n)}e=0$ for every $n\geq0$.
Require also $e-\mathbf1\otimes S\in V\otimes_D B$; then
$e_{(0)}$ preserves $V\otimes_D B$ and squares to zero there.
Morphisms preserve every mode, translation, grading, filtration,
inclusion, vacuum and the two marked states $S,e$.
They preserve $d,w$ exactly. They need not preserve the names of
generators; derivatives and composites may occur in their images.
The Miura system just constructed is denoted by $\mathcal F_R$.
Its state modules are direct sums of finite PBW words with arbitrary
finite jets.
\end{definition}

To allow resolutions, give these state modules a second, resolution
grading $r$, retaining the ghost grading $g$. The tensor signs use
total parity $r+g$ modulo two. A resolution has a $D$-linear
derivation $\partial$ of bidegree $(1,0)$ and weight zero,
commuting with translation and preserving the filtration, with
\[
 \partial^2=0,\qquad \partial S=\partial e=0,
 \qquad \partial d+d\partial=0.
\]
Here $S,e$ have bidegree $(0,1)$, $w$ has bidegree $(0,2)$,
and $d$ has bidegree $(0,1)$. The original system has resolution
degree zero and $\partial=0$. The curved derivation $d$ is not
used to define resolution cohomology.

Invert the morphisms that are quasi-isomorphisms for $\partial$
on each of the three state complexes. Equivalences and coherent
homotopies below are taken in this localization. Use representatives
whose underlying state complexes are K-flat when forming derived
scalar extensions. In the equivariant category the morphisms and
homotopies intertwine the specified parameter action. These choices
allow acyclic resolution states and invertible filtered changes of
fields while retaining the underlying derived fibre representation.
They impose no field/bar comparison as an entry hypothesis.

\begin{definition}\label{rpd:def:descent}
A \emph{coarse descent} of the signed system $\mathcal F_R$ is a
field system $\mathcal D$ over $A$ and an equivalence
\[
 \theta:R\otimes_A^{\mathbf L}\mathcal D\simeq\mathcal F_R
\]
in this category, identifying the canonical semilinear action on the
scalar factor with $\sigma$. The induced comparison over
$R\otimes_A R$ must satisfy the parameter \v{C}ech cocycle, with
its coherent homotopies for resolved systems.
\end{definition}

\begin{theorem}[Failure of coarse-parameter descent]
\label{rpd:thm:descent}
The signed universal Miura $W_4$ field system has no descent from
$R=\mathbb C[a]$ to $A=\mathbb C[a^2]$ in the preceding category.
The same obstruction persists after equivariant filtered changes of
states and after equivariant resolution equivalences. The system descends after
inverting $a^2$, and it defines a field system on the quotient stack
$[\operatorname{Spec}R/\mu_2]$ without inverting $a$.
\end{theorem}
\begin{proof}
The parameter \v{C}ech relation is
\[
 R\otimes_A R=\mathbb C[a,b]/(a^2-b^2).
\]
Its diagonal $b=a$ and its reflection graph $b=-a$ meet at
$a=b=0$. A descent comparison restricts to the identity on the
diagonal and to the specified action on the reflection graph.
On $\gamma_3$ these restrictions at their intersection would be
$\gamma_3$ and $-\gamma_3$. Their difference is
$2\gamma_3\ne0$, by the free PBW basis in
Lemma~\ref{rpd:lem:polynomial}.

The same obstruction can be expressed on a derived fibre. For a
pulled-back $A$-complex $C$, associativity of derived scalar extension
gives an equivariant identification
\[
 \mathbb C\otimes_R^{\mathbf L}
       (R\otimes_A^{\mathbf L}C)
 \simeq \mathbb C\otimes_A^{\mathbf L}C.
\]
Compute the right side using a K-flat $A$-representative of $C$.
The canonical involution changes only the $R$-factor before restriction;
after restriction it is the identity on every term. This remains true
if $C$ is not flat over $A$. The fibre of
$B_R$ has a nonzero $(-1)$-eigenclass $\gamma_3$ in resolution
cohomology $H^0_\partial$. An equivariant equivalence induces an
isomorphism of these representations. No change of frame, homotopy,
or acyclic resolution can turn that representation into the identity
representation. This also proves the assertion for homotopy-coherent
descent, rather than only strict descent.

After inverting $a^2$, the two graphs are disjoint and the cover is a
$\mu_2$-torsor. On a PBW word $v$ of sign $(-1)^\varepsilon$, the
comparison $p_1^*\mathcal F_R\longrightarrow p_2^*\mathcal F_R$,
where $p_1,p_2$ are the two projections from the parameter fibre
product, is
\[
 \Theta_{a,b}(v)=(b/a)^\varepsilon v.
\]
Its restrictions are the identity and the involution. Since
$(b/a)^2=1$, coefficient covariance proves preservation of all modes
and of the coupling. On three parameter lifts its cocycle is the
identity $(h/b)^\varepsilon(b/a)^\varepsilon=(h/a)^\varepsilon$.
Here $h$ is the third parameter coordinate. An explicit descended
system is obtained by taking the invariant field algebras over
$A[t^{-1}]$. Positive PBW words remain a basis, while a negative word
$v$ is replaced by $av$; their pullback comparison is invertible
because $a$ is invertible. Its exact coupling is
\[
 e'=\mathbf1\otimes S+U_2\otimes\gamma_2+U_4\otimes\gamma_4
       +t^{-1}(aU_3)\otimes(a\gamma_3).
\]
Its pullback is $e$. All mode, curvature and coupling identities follow
by the injective faithfully flat scalar extension, and the invariant
subalgebras carry the restricted translation, inclusion and vacuum.
Without
inverting $a$, the action itself still supplies descent along the
atlas of the quotient stack. That action groupoid is different from
the \v{C}ech relation of the coarse parameter space at the branch.
\qedhere
\end{proof}

This theorem does not exclude a form obtained by forgetting the signed
action or choosing a different lift of parameter reflection. It also
does not assert invariance under a localization that retains only the
cohomology of the coupled differential $Q$.

\section{The comparison map and the coupling}

Let $\mathcal P_+$ and $\mathcal P_-$ be the PBW words on which the
constant sign part of $\sigma$ is respectively $+1$ and $-1$.
Derivatives have the sign of the state they differentiate. Normal
words are exact eigenvectors, because $\sigma$ is a vertex
homomorphism.

\begin{proposition}[Invariant comparison]
\label{rpd:prop:invariants}
For $B_R$, the invariant module is
\[
 B_R^\sigma=\bigoplus_{v\in\mathcal P_+}A v
       \;\oplus\bigoplus_{v\in\mathcal P_-}A(av).
\]
The natural field homomorphism
\[
 \lambda:R\otimes_A B_R^\sigma\longrightarrow B_R,
 \qquad r\otimes x\longmapsto rx,
\]
acts by $1$ on a positive word and by $a$ on a negative word. It is
injective, and
\[
 H^0\operatorname{Cone}(\lambda)
   =\bigoplus_{v\in\mathcal P_-}R/(a)\,v,
 \qquad H^{-1}\operatorname{Cone}(\lambda)=0.
\]
On the derived fibre at $a=0$, every negative word contributes one
copy of $\mathbb C$ in each of resolution degrees $-1$ and $0$.
\end{proposition}
\begin{proof}
An invariant coefficient of a positive word is an even polynomial in
$a$, hence lies in $A$. An invariant coefficient of a negative word is
an odd polynomial, hence lies in $aA$. This proves the decomposition
and the formula for $\lambda$. Multiplication by $a$ on $R$ is
injective with cokernel $R/(a)$. Since $R$ is free of rank two over
$A$, the source tensor product is already the derived tensor product.
On a negative word, tensoring the two-term cone with $R/(a)$ makes
its differential zero,
which proves the assertion about the derived fibre.

The cone is independent of the PBW frame used to compute it.
Averaging by $(1+\sigma)/2$ makes invariants exact on the underlying
$A$-complexes, and $R$ is flat over $A$. Naturality of $\lambda$
therefore identifies its cones under equivariant resolution
quasi-isomorphisms. A system pulled back from $A$ has an invertible
$\lambda$, since $(R\otimes_A C)^\sigma=C$ with its canonical
action. Thus the displayed cone is also an intrinsic necessary
obstruction to coarse descent in the specified category.

The same argument applies to $V_R$ and $B_R^\circ$ with their own
PBW bases.\qedhere
\end{proof}

\begin{proposition}[Tensor comparison]\label{rpd:prop:tensor}

The natural comparison
\[
 \mu:V_R^\sigma\otimes_A B_R^\sigma
       \longrightarrow (V_R\otimes_R B_R)^\sigma
\]
has matrix $\operatorname{diag}(1,1,1,t)$ on the sign sectors
$(+,+),(+,-),(-,+),(-,-)$, in that order.
Its cokernel is a copy of $A/(t)$ for each negative--negative pair of
PBW words. In particular, the invariant coupling term
$U_3\otimes\gamma_3$ is not in its image.

The same comparison with $B_R^\circ$ has the same four-sector form.
Consequently the full exact coupling $e$ does not lift to
$V_R^\sigma\otimes_A(B_R^\circ)^\sigma$.
\end{proposition}
\begin{proof}

For a positive--negative or negative--positive pair, the factor $a$
already occurs in the invariant basis of the target. On a
negative--negative pair the source basis is $(av)\otimes(ab)$ and
the target basis is $v\otimes b$. The map is multiplication by
$a^2=t$. PBW independence prevents other pairs from supplying the
missing coefficient.\qedhere
\end{proof}

The map $\lambda$ preserves every mode. Its cone measures the states
that its source fails to recover.

\begin{example}\label{rpd:ex:collision}
For a concrete nonlinear operation put $\chi=\gamma_4$ and
$p=\chi_{(-1)}\chi_{(-2)}\mathbf1$. The dual relation gives
\[
 \gamma_{3,(0)}\gamma_3=\tfrac34p,\qquad
 (a\gamma_3)_{(0)}(a\gamma_3)=\tfrac34t\,p.
\]
The state $p$ is nonzero: its two odd creation pairs are distinct.
Thus invariants do not remove the noncommutative relation, but they
multiply its coefficient by the ramified parameter. The state $a\gamma_3$
is a free $A$-basis element of the invariant module: its class modulo
$t$ is nonzero there, even though $\lambda$ sends that class to zero
in the target fibre at $a=0$.
\end{example}

The full coupled algebra $(V_R\otimes B_R,Q)^\sigma$ does exist over
$A$. Proposition~\ref{rpd:prop:tensor} says that it cannot be
identified with the tensor product of separately invariant factors
by the natural comparison. Existence of an invariant total coupling
is not a descent theorem for its two factors.

\section{Completion, curvature, and changes of fields}

Put $\widehat A=\mathbb C[[t]]$ and $\widehat R=\mathbb C[[a]]$.
Complete the field modules in the parameter, not in PBW word length:
\[
 \widehat B=\varprojlim_n B_R/a^nB_R.
\]
An element is a family of PBW coefficients tending to zero
$a$-adically: modulo each $a^n$, only finitely many coefficients
are nonzero. This is a restricted family, not an unrestricted product
of state spaces. Locality is required at each finite parameter
precision. Completed tensor products here mean inverse limits of the
corresponding tensor products modulo $t^m$, or equivalently along the
cofinal tower $a^{2m}$. The gradings extend by continuous projections
to the homogeneous components; only adically restricted sums of such
components are allowed. These are topological field algebras, with
the vertex identities and locality imposed modulo every parameter power.

Invariants commute with this inverse limit, since they are the
equalizer of $\sigma$ and the identity. Exactness of averaging
identifies $\widehat B^{\sigma}$ with the $t$-adic completion of
$B_R^\sigma$. The completed comparisons are
\[
 \widehat\lambda:\widehat R\widehat\otimes_{\widehat A}
       \widehat B^\sigma\longrightarrow\widehat B,
 \qquad
 \widehat\mu:\widehat V^\sigma\widehat\otimes_{\widehat A}
       \widehat B^\sigma
 \longrightarrow(\widehat V\widehat\otimes_{\widehat R}\widehat B)^\sigma.
\]

\begin{proposition}\label{rpd:prop:completion}
The completed comparison $\widehat\lambda$ is still injective and has
cokernel $\bigoplus_{v\in\mathcal P_-}\mathbb C v$. The completed
tensor comparison still acts by $t$ on every negative--negative
sector. Neither completion restores descent at the branch.
\end{proposition}
\begin{proof}
On each negative PBW coefficient the first map is multiplication by
$a$ on $\mathbb C[[a]]$. Division by $a$ of a restricted family all
of whose coefficients are divisible by $a$ remains a restricted
family. Thus its image consists exactly of negative families with
zero constant coefficient. Only finitely many constant coefficients
can be nonzero, proving the stated cokernel and injectivity.

At precision $a^{2m}$, multiplication by $a$ has kernel generated by
$a^{2m-1}$ and cokernel $\mathbb C$. The kernel maps to zero at the
previous even precision. The finite kernels therefore do not give a
kernel in the limit. The tensor assertion follows coefficientwise
from multiplication by $a^2=t$ on the negative--negative sector.
This direct description specifies the passage to the inverse limit.

For resolved completed systems, $\widehat R$ is free of rank two
over $\widehat A$, so its completed scalar extension is the ordinary
finite-free scalar extension on complete state modules. The derived
branch fibre of a pulled-back system again has trivial involution.
Multiplication by $a$ is injective on $\widehat B$, and
$\widehat B/a\widehat B$ is the direct sum of the residual PBW words.
The two-term free resolution of $\mathbb C$ over $\widehat R$
therefore identifies this quotient with the derived fibre of
$\widehat B$. Its class $\gamma_3$ remains nonzero and negative.
The representation argument of Theorem~\ref{rpd:thm:descent}
excludes descent after every permitted equivariant resolution.\qedhere
\end{proof}

\begin{example}\label{rpd:ex:scalar}
The scalar term is retained at the obstructed fibre. Put $L=-U_2$
as in the Miura normalization. At $a=0$ the
central charge is $3$, so $L_{(3)}L=\tfrac32\mathbf1$. A unital vertex
augmentation to the scalars would send its left side to zero and its
right side to $3/2$, and therefore cannot exist. The marked vacuum is
nevertheless nonzero and is fixed by the involution. The curvature
witness specializes to
\[
 d^2\chi=w_{(0)}\chi
   =\tfrac32\gamma_{2,(-4)}\chi_{(-1)}\chi_{(-2)}\mathbf1\ne0.
\]
Replacing $w$ by $2w$, while keeping $d$ fixed, gives the nonzero defect
$-\tfrac32\gamma_{2,(-4)}\chi_{(-1)}\chi_{(-2)}\mathbf1$.
These identities specialize Proposition~\ref{w4d:prop:curvature};
the three creation pairs are distinct, so the displayed state is
nonzero by Lemma~\ref{rpd:lem:polynomial}.
\end{example}

A primary change of fields does not remove the obstruction. The
primary denominator $37-300a^2$ of Proposition~\ref{w4d:prop:primary}
equals $37$ at the branch. The change is therefore defined on the
open set where that denominator is invertible, and on the completed
branch neighbourhood. Its odd dual state is
\[
 \eta_3=\gamma_3-\tfrac a2T\chi,
 \qquad \sigma(\eta_3)=-\eta_3,
 \qquad \eta_3|_{a=0}=\gamma_3.
\]
Any invertible equivariant change of states conjugates the action on
the fibre and preserves its nontrivial negative eigenspace. In
transporting the coupling one must retain both its nonlinear
composite term and its total derivative, as in the exact primary
coupling formula; a coupling modulo translation is not the marked
coupling used here.

The three loci relevant to these arguments are distinct:
\[
 \begin{array}{c|c|c}
  &a^2&c\\ \hline
 \text{ramification}&0&3\\
 \text{vanishing curvature}&1/20&0\\
 \text{primary denominator zero}&37/300&-22/5.
 \end{array}
\]
Every nonnegative product $S_{(n)}S$ vanishes on the middle locus, as
follows from the common factor $20a^2-1$ in the scalar-dual row.
The Miura tables have no denominator at the third locus, so the
universal polynomial specialization exists there as well.

\section{The geometric comparison over the parameter stack}

The preceding construction concerns field state modules on a fixed
coordinate chart. Question~\ref{w4d:qn:bar} asks for the additional
comparison with completed ordered collision coefficients. Such a
comparison must preserve cut convolution and the diagonal tensor,
their separate units $K,J$, and the supported operations. The
nonconsecutive stratum $13|2$, normal derivatives and spectator states
are part of those operations.

\begin{corollary}\label{rpd:cor:geometric}
Suppose a geometric enhancement of $\mathcal F_R$ has an equivariant
state-extraction functor that preserves equivalences and commutes
with derived pullback along $A\to R$. If its extraction is equivalent
to the signed system $\mathcal F_R$, the enhancement has no coarse
descent whose extracted comparison is a field-system comparison of
Definition~\ref{rpd:def:descent}.
\end{corollary}
\begin{proof}
Applying extraction to such a descent would produce the descent of
$\mathcal F_R$ excluded by Theorem~\ref{rpd:thm:descent}.
\end{proof}

\begin{question}\label{rpd:qn:geometric}
Does the field/bar comparison of Question~\ref{w4d:qn:bar} admit an
equivariant construction over $[\operatorname{Spec}R/\mu_2]$, with
independently defined state extraction and coherent parameter
base-change maps? Can its supported collision maps, its unit and
counit, and both geometric product comparisons be constructed there,
retaining the exact coupling and the nonempty vacuum?
\end{question}

The quotient-stack field system supplies the parameter action required
by this question. The coefficient maps and a cooperation compatible
with the chiral product remain to be constructed.

Translation is already null-homotopic for $Q$ over $R$.
The degree-minus-one vertex derivation $h$ of
Proposition~\ref{w4d:prop:cartan} is determined by
\[
 h\gamma_2=-\mathbf1,\qquad h\gamma_3=h\gamma_4=0,
\]
and is extended by zero on $V_R$. On a PBW word it removes an
occurrence of the creation mode $\gamma_{2,(-1)}$ with coefficient
$(-1)^{r+1}$, where $r$ is the number of preceding odd factors,
and sums over such occurrences. Hence it preserves $B_R$.
The identity $[Q,h]=T_V+T_B$ holds after extension to $\mathbb C(a)$,
and therefore on $V_R\otimes_R B_R$, which is free over $R$.
Moreover $h$ commutes with $\sigma$: the derivations $h$ and
$\sigma h\sigma$ have the same values on the generators.
Question~\ref{w4d:qn:topologization} asks for a geometric realization
of this Cartan homotopy on $\mathbb R\times X$, preserving the full
coupling and the collision products.
